% Verbatim copy of /Users/alingalatan/workspaces/grothe/handoff/manuscript/legacy/ch25-parts/70-canonical-ellipse.tex (file dated 2026-10-25). Only this line was added.
\subsection{Canonical response geometry, a closed-form ellipse, and a retained phase defect}
\label{sec:canonical-response-ellipse}

\paragraph{Status and scope.}
[P] denotes an analytic proof and [C] a directed exact calculation.
This section reorganizes the original-response route, strengthens its
pointwise inequality, and proves a new region. All statements allow every
visible sector, every complex phase, every finite dimension, and arbitrary
sharing of the remote Gaussian coordinates. No moment inequality,
stationarity, raw Conjecture K, or fresh Gaussian buffer is assumed.
The full complement and a new global Grothendieck endpoint remain open.

Fix $\epsilon=1891/2500$ and $p=1/\epsilon$. For disk-valued responses put
\[
 h=(f_++f_-)/2,\quad k=(f_+-f_-)/2,\quad |f_\pm|\le1,
 \quad Ph=tZ,\quad F=\epsilon Qk,
\]
where $P=P_{1,0}$ and $Q=P_{2,1}$. The visible proper complex Gaussian
$Z$ is independent of the remote Gaussian space $X$. Write $T=|Z|^2$ and
\[
 F=q+(T-1)r_A+Z^2r_B/\sqrt2+ZY+\bar ZW+\nu H(Z),
 \qquad H(Z)=Z(T-2)/\sqrt2.
\]
The remote bidegrees are $(2,1),(1,0),(0,1),(1,1),(2,0)$, respectively.
Set $e=\|q\|^2$, $A_2=\|r_A\|^2$, $B_2=\|r_B\|^2$, and
$E_2=\|Y\|^2+\|W\|^2$. The objective is exactly
\[
 \mathcal N=t^2+p\|F\|^2-\|Pk\|^2-\epsilon\|Qh\|^2.
\]
In particular, $q=\mathbb E[F\mid X]$ is the actual, unnormalized remote
projection. It is not an arbitrary unit cubic inserted into this formula.

\subsubsection{The canonical target is a convex response body [P, synthesis]}
In fixed finite Gaussian dimension set $\mathcal E=\operatorname{ran}P
\oplus\operatorname{ran}Q$. Let $\mathcal C_n$ be the set of pairs
$(x,y)=(P h+Qh,Pk+Qk)$ obtained from $|h\pm k|\le1$.
It is a compact convex set in the underlying finite-dimensional real
space: it is the continuous finite-moment image of two weak-* compact
$L^\infty$ disks. For every $u,v\in\mathcal E$ its support function is
\[
 \sup_{(x,y)\in\mathcal C_n}2\Re\{\langle u,x\rangle+\langle v,y\rangle\}
       =\mathbb E|u+v|+\mathbb E|u-v|.
\]
Indeed the left side separates into the two disk support functions for
$f_+$ and $f_-$. Conversely, all these support inequalities characterize
$\mathcal C_n$, by finite-dimensional separation. Thus the exact norm
problem is
\[
 \sup_{(x,y)\in\mathcal C_n}
 \{\|Px\|^2-\|Py\|^2-\epsilon\|Qx\|^2+\epsilon\|Qy\|^2\}.
\]
The required theorem bounds this quantity uniformly in $n$. Raw K is one
sufficient relaxation after dropping response constraints and negative
terms. The response-body formulation retains them; it does not require
that raw K hold for every abstract cubic. This is a synthesis of the
existing exact response duals, not an additional endpoint theorem.

\subsubsection{New regions [P,C]}
For every such response,
\begin{equation}\label{eq:cre-global-quadratic}
\boxed{\begin{aligned}
 \mathcal N\le{}&\frac{1195}{2000}+\frac{1961}{10000}t^2
       +\frac{257}{500}e
       -\frac{A_2+E_2}{1250}-\frac{101}{250}B_2
       -\frac{33}{100}|\nu|^2.
\end{aligned}}
\end{equation}
Thus the following ellipse is proved:
\begin{equation}\label{eq:cre-ellipse}
 \boxed{\quad e+\frac{191}{500}t^2\le\frac9{40}
 \quad\Longrightarrow\quad
 \mathcal N\le\frac{14263}{20000}=0.71315<\frac{500}{701}.\quad}
\end{equation}
For example it allows $e\le0.225$ at $t=0$, $e\le0.1295$ at $t=0.5$,
and $e\le0.03782$ at $t=0.7$. This does not bound all possible responses
at those $t$ values, since they can have larger remote variance.

Combining this ellipse with the preceding anisotropic certificate gives
\begin{equation}\label{eq:cre-union}
 \boxed{\quad 0\le t\le\frac{851}{1000},\qquad e\le\frac{37}{2500}
 \quad\Longrightarrow\quad
 \mathcal N\le\frac{89157}{125000}=0.713256<\frac{500}{701}.\quad}
\end{equation}
This extends the earlier interval $[0.667,0.851]$ all the way to zero.
Indeed the ellipse covers $0\le t\le37/50=0.74$ at $e\le37/2500$,
and the anisotropic certificate covers $[0.667,0.851]$. Both are
recomputed by the portable checker; no moment-dependent radius estimate
is needed for this union.

\subsubsection{The stronger two-disk inequality is one square [P]}
For every $L\in\mathbb C$, $\lambda>0$, and $|h\pm k|\le1$,
\begin{equation}\label{eq:cre-full-disk}
 \boxed{\quad
 2\Re(\bar Lh)+\lambda|k|^2+
 \frac{[\Re(\bar Lk)]^2+[\Im(\bar Lh)]^2}{\lambda}
 \le \lambda+\frac{|L|^2}{\lambda}.
 \quad}
\end{equation}
The previous two-disk inequality omitted the last, nonnegative term.
To prove the stronger form, rotate $L$ to $l=|L|$ and first take
$f_\pm$ on the unit circle. Write
$h=e^{i\phi}\cos\psi$, $k=ie^{i\phi}\sin\psi$, and set
$a_0=\cos\phi$, $c_0=\cos\psi$. The gap in
\eqref{eq:cre-full-disk}, multiplied by $\lambda$, is exactly
\[
 \lambda^2+l^2-2\lambda l a_0c_0-
 \lambda^2(1-c_0^2)-l^2(1-a_0^2)
       =(\lambda c_0-la_0)^2.
\]
For disk-valued $f_\pm$, express each as a convex combination of two
antipodal circle points and use convexity of the entire left side.
The new coefficient is sharp with the other terms fixed: take
$0<a_0<\min\{1,\lambda/l\}$ and $c_0=(l/\lambda)a_0$ when $l>0$.
Then the gap is zero while $\Im h\ne0$.

This proof needs no Gaussian structure. That structure enters only when
the resulting weighted norms are projected onto Hermite components.

\subsubsection{Exactly which quadratic majorants are possible [P]}
The pointwise geometry admits a complete elementary description.
For $l\ge0$ and $A,B\ge0$, let
\[
 \Psi_l(A)=
 \begin{cases}2l,&0\le A\le l,\\ A+l^2/A,&A>l.\end{cases}
\]
When $l=0$, interpret $\Psi_0(0)=0$. Then
\begin{equation}\label{eq:cre-exact-envelope}
 \sup_{|h\pm k|\le1}
 \{2l\Re h+A(\Im k)^2+B(\Re k)^2\}
       =\max\{B,\Psi_l(A)\}.
\end{equation}
Convexity reduces to circle-valued $f_\pm$. In the preceding coordinates
the expression is
\[
 B(1-c_0^2)+2l a_0c_0-(B-A)a_0^2(1-c_0^2).
\]
One may take $a_0,c_0\in[0,1]$. If $B\le A$, maximizing in $a_0$
gives $a_0=1$, followed by $\Psi_l(A)$. If $B>A$, the interior
maximizer is $a_0=lc_0/((B-A)(1-c_0^2))$ while that value is at most
one. Substitution, with $u=c_0^2$, gives
\[
 B(1-u)+\frac{l^2u}{(B-A)(1-u)}.
\]
This is convex in $u$, so its maximum on the interior branch is at an
endpoint: either $B$, or a boundary value already bounded by
$\Psi_l(A)$. The lower bounds are attained by $h=0,k=1$ and by the
$a_0=1$ family, respectively. The case $l=0$ follows directly.

Consequently, for a prescribed constant $C\ge2l$, the largest
reflection-invariant nonnegative coefficients on the two components of
$k$ are
\[
 B=C,\qquad A=\frac{C+\sqrt{C^2-4l^2}}2.
\]
Equivalently, they are $A=\lambda$, $B=\lambda+l^2/\lambda$ with
$\lambda\ge l$. This characterizes this particular quadratic family;
it is not a no-go theorem for nonlinear response inequalities. The
retained $h$ term in \eqref{eq:cre-full-disk} is additional information.

\subsubsection{A two-parameter radial certificate with analytic profile elimination [P]}
Choose any positive radial weight $\lambda(T)$ with the following
moments finite:
\[
 J_j=\mathbb E\frac{T^j}{\lambda(T)}\quad(0\le j\le3),
 \qquad W_0=\mathbb E\lambda(T).
\]
Define
\[
\begin{split}
 K_A&=J_2-2J_1+J_0-\frac{(J_1-J_0)^2}{1+J_0},\\
 K_B&=J_2/2,\\
 K_\nu&=\frac{J_3-4J_2+4J_1}{2}
       -\frac{(J_2-2J_1)^2}{2(1+J_1)}.
\end{split}
\]
Weighted Bessel with the first-chaos penalty gives
\begin{equation}\label{eq:cre-isotropic-bessel}
 \mathbb E\lambda|k|^2+\|Pk\|^2\ge p^2
 \left\{\frac e{J_0}+\frac{A_2}{K_A}+\frac{B_2}{K_B}
             +\frac{E_2}{J_1}+\frac{|\nu|^2}{K_\nu}\right\}.
\end{equation}
Here all denominators must be positive. The $A$ block couples the two
visible factors $T-1$ and $1$ on the same remote first-chaos direction;
the diagonal block couples $Z(T-2)/\sqrt2$ with $Z$. Adding $\|Pk\|^2$
gives exactly the two displayed rank-one inverse corrections. The other
remote Hermite blocks are orthogonal, so this proof is dimension-free.

For a real auxiliary profile $L_t=Z(a+bT)$, put
$\mathcal B_t=bZ(T-2)$. The exact completion
\[
 -\epsilon\|Qh\|^2=2\Re\langle\mathcal B_t,h\rangle+2pb^2
                 -\epsilon\|Qh+p\mathcal B_t\|^2
\]
and \eqref{eq:cre-full-disk} imply
\begin{align*}
 \mathcal N\le{}&t^2-2(a+2b)t+W_0+a^2J_1+2abJ_2+b^2(J_3+2p)\\
 &+(p-p^2/J_0)e-(p^2/K_A-p)A_2-(p^2/K_B-p)B_2\\
 &-(p^2/J_1-p)E_2-(p^2/K_\nu-p)|\nu|^2-\mathcal R_t,
\end{align*}
where the nonnegative, compatible remainder is
\begin{equation}\label{eq:cre-defect}
 \mathcal R_t=
 \mathbb E\frac{[\Re(\overline{L_t}k)]^2+
                    [\Im(\overline{L_t}h)]^2}{\lambda(T)}
       +\epsilon\|Qh+p\mathcal B_t\|^2.
\end{equation}
Only the isotropic $k$ norm was used in
\eqref{eq:cre-isotropic-bessel}; neither angular term has been charged
there. The $Pk$ square has been used and must not be subtracted again.

The profile parameters need not be searched over. Let
\[
 H=\begin{pmatrix}J_1&J_2\\J_2&J_3+2p\end{pmatrix},\qquad d=\binom12.
\]
This matrix is positive definite, by weighted Cauchy--Schwarz and $p>0$.
Choose $(a,b)^{\mathsf T}=tH^{-1}d$. Then the polynomial part of the bound is
\begin{equation}\label{eq:cre-master}
\boxed{\begin{aligned}
 \mathcal N\le{}&W_0+(1-d^{\mathsf T}H^{-1}d)t^2+(p-p^2/J_0)e\\
 &-(p^2/K_A-p)A_2-(p^2/K_B-p)B_2\\
 &-(p^2/J_1-p)E_2-(p^2/K_\nu-p)|\nu|^2-\mathcal R_t.
\end{aligned}}
\end{equation}
This is a reusable certificate, simultaneously for every $t$. It avoids
row-by-row profile optimization. In the present example $a/t$ and $b/t$
are positive, but no sign assumption on them was needed for the argument.

\subsubsection{Four moments from one closed form [P,C]}
For $s,d_0>0$ take $\lambda(T)=\sqrt{s+d_0T}$ and set $z=s/d_0$.
The complete moment calculation is
\begin{align*}
 J_0&=\frac{\sqrt\pi}{\sqrt{d_0}}e^z\operatorname{erfc}(\sqrt z),\\
 J_1&=(1/2-z)J_0+\frac{\sqrt s}{d_0},\\
 J_{j+1}&=(j+1/2-z)J_j+jzJ_{j-1}\quad(j=1,2),\\
 W_0&=\sqrt s+\frac{d_0}2J_0.
\end{align*}
The first line is the substitution $u=\sqrt{T+z}$ in the defining
integral. Integrating the derivative of
$e^{-T}T^j\sqrt{s+d_0T}$ proves the recurrence, with its boundary term
retained at $j=0$. Thus there is no quadrature or tail cutoff in this new
certificate.

Set $s=43/2000$ and $d_0=417/1000$. Exact intervals prove
\[
\begin{array}{c|c}
 W_0&[0.5974205484163568,\ 0.5974205484163569]\\
 1-d^{\mathsf T}H^{-1}d&[0.1960124339463171,\ 0.1960124339463172]\\
 p-p^2/J_0&[0.5136504471928151,\ 0.5136504471928152]
\end{array}
\]
and the penalties for $(A_2,B_2,E_2,|\nu|^2)$ exceed, respectively,
\[
 1/1250,\qquad101/250,\qquad1/1250,\qquad33/100.
\]
This proves \eqref{eq:cre-global-quadratic}; the same rational inequality
can retain $-\mathcal R_t$. Since
$1961/10000<(257/500)(191/500)$, the ellipse cap is exactly
\[
 \frac{1195}{2000}+\frac{257}{500}\frac9{40}
       =\frac{14263}{20000}<\frac{500}{701}.
\]

The standalone checker uses integer intervals at scale $10^{40}$.
It bounds $\pi$ using Machin's identity and alternating arctangent series,
$e^z$ using a positive series with an explicit remainder, and
$\operatorname{erf}(\sqrt z)$ using its alternating series. It then
performs the recurrence and every matrix comparison with outward rounding.
The union theorem additionally replays the pinned anisotropic certificate.

\subsubsection{An extra term in the preceding anisotropic theorem [P]}
For that theorem's exact auxiliary profile and weight,
\[
 L_*=Z\left(\frac{4541}{10000}+\frac{3049}{20000}T\right),\qquad
 \lambda_*^2=|L_*|^2+\frac{167}{50000}+\frac{11}{500}T,
\]
define
\[
 \mathcal A_h=\mathbb E\frac{[\Im(\overline{L_*}h)]^2}{\lambda_*}.
\]
Every bound in its derivation can retain $-\mathcal A_h$, together with
its already available $-\epsilon\|Qh+p\mathcal B_*\|^2$.
The previous $k$-parallel term was used in its real-linear inverse Gram
matrix; it cannot also be subtracted a second time.
In particular a central-row violator must satisfy
\[
 \mathcal E_{\rm adj}\ge
 \frac{500/701-t^2+1.518t-1.272881}{0.54}
       +\frac{\mathcal A_h+\epsilon\|Qh+p\mathcal B_*\|^2}{0.54},
\]
where $\mathcal E_{\rm adj}$ is precisely the preceding anisotropic
sector-adjusted variance. This is an additional necessary constraint,
not a proof that the surviving set is empty.

\subsubsection{Canonical bookkeeping and the remaining problem}
The reusable chain now has three stages: a pointwise response inequality,
a finite weighted-Gram calculation by remote Hermite type, and a scalar
or matrix certificate for its coefficients. The last two stages do not
depend on the ambient dimension. Raw-K families provide different
certificates on restricted cubic laws; they are not interchangeable with
conditions on actual projections of bounded responses.

The preceding cumulant/support theorem controls another region and keeps
its own moment dependency. The present region is moment-free. One may
take the minimum of the valid upper bounds, but may not add their separate
penalties if they already spent the same $Pk$ or $Qh$ square.
The new ellipse, the strengthened anisotropic region, and the support
criteria do not yet cover all actual responses. The full same-$\epsilon$
row cover and the existing reduction/transfer chain remain necessary;
the certified global interval is unchanged.
